\uuid{mhpU}
\titre{Permis de conduire en deux tentatives}
\auteur{worms}
\organisation{exo7}
\licence{CC BY-NC-SA 4.0}
\datecreate{2017-08-28}
\niveau{L2}
\module{Probabilité et statistique}
\chapitre{Probabilité discrète}
\sousChapitre{Probabilité et dénombrement}
\difficulte{}
\theme{Probabilités, événements}
\exo7qcmid{worms.worms.005}
\serie{ordonnee}

\contenu{
\texte{
En France, on considère la population ${\cal P}$ des candidats au permis de conduire qui essaient de l'obtenir une fois, puis une seconde fois si la première tentative échoue. Parmi eux, un candidat sur trois l'obtient du premier coup, et parmi ceux qui ne l'ont pas eu du premier coup, 30\% d'entre eux l'obtiennent à la seconde tentative.
On considère l'expérience aléatoire consistant à sélectionner au hasard une personne issue de cette population ${\cal P}$.
}

\begin{question}
\qclassification{260.01}
Dans ce cadre, le(s)quel(s) des 2 espaces $\Omega$ ci-dessous peu(ven)t être considéré(s) pour cette expérience ?
\begin{answers}
\bad{Seulement $\Omega=\{$ permis obtenu \ , \ permis non obtenu  $\}$.}
\bad{Seulement $\Omega=\{$\ il y a eu une tentative \ , \ il y a eu deux tentatives $\}$.}
\bad{Aucun des deux n'est un univers adéquat.}
\good{Les deux peuvent convenir.}
\end{answers}
\begin{explanations}
Les deux peuvent en effet convenir, mais sont des univers pauvres. Par exemple, "le permis a été obtenu à la première tentative" est un événement pour l'un mais pas pour l'autre, et c'est l'inverse pour "le permis a été obtenu"...
\end{explanations}
\end{question}

\begin{question}
\qclassification{260.01}
On suppose désormais qu'un espace $\Omega$ convenable a été choisi (mais on ne le détaille pas ici ; il permet en tout cas de définir les événements adéquats des questions suivantes). La probabilité d'obtenir le permis au plus tard à la seconde tentative vaut :
\begin{answers}
\bad{$1/2$}
\bad{$2/3$}
\good{Environ $53\%$}
\bad{Environ $23\%$}
\end{answers}
\end{question}

\begin{question}
\qclassification{260.01}
Peut-on définir les événements $A=$"la seconde tentative a échoué sachant que la première a échoué"
et $B=$"la première tentative a échoué et la seconde a réussi" ?
\begin{answers}
\bad{Oui pour $A$, oui pour $B$.}
\bad{Oui pour $A$, non pour $B$.}
\good{Non pour $A$, oui pour $B$.}
\bad{Non pour $A$, non pour $B$.}
\end{answers}
\begin{explanations}
A n'est pas un événement ! En effet, à la question "Est-ce que la seconde tentative a échoué sachant que la première a échoué ?", on ne peut pas répondre par oui ou par non, car la question n'a aucun sens. Ce n'est pas le cas de B, qui définit bien un événement.
\end{explanations}
\end{question}

\begin{question}
\qclassification{260.01}
\qkeeporder
Que vaut la probabilité d'avoir tenté l'épreuve une seconde fois sachant qu'on a obtenu le permis au final ?
\begin{answers}
\bad{Zéro.}
\good{$0,375$}
\bad{$0,1875$}
\bad{Une autre valeur que les réponses précédentes.}
\bad{La question n'a pas de sens.}
\end{answers}
\end{question}

\begin{question}
\qclassification{260.01}
Les événements $R=$"le permis est obtenu à l'issue de l'expérience" et $T=$"Une deuxième tentative a eu lieu" sont-ils :
\begin{answers}
\bad{Indépendants.}
\bad{Incompatibles.}
\bad{Indépendants et incompatibles.}
\good{Ni l'un ni l'autre.}
\end{answers}
\end{question}
}
